\section{Testable predictions and experimental/data-analysis protocols}
\label{sec:tests}

The appropriate empirical posture for HPA--$\Omega$ is to specify testable interfaces: statistical biases, scaling laws, and protocol-level certificates that can be measured across systems.
This section gives three minimal, falsifiable templates aligned with the AEC definition.

\subsection{Audit posture: bounded-complexity protocols}
\label{sec:test_audit_posture}
Each test below depends on explicit protocol choices: how a readout variable is mapped to a phase coordinate (uniformization), how ratios and tolerances are inferred, which window length $N$ and step size $\Delta t$ are used, which smoothing operator is applied, and which surrogate/baseline family is used for comparison.
To avoid post-hoc freedom, these choices should be fixed \emph{a priori} or selected by a deterministic bounded-complexity rule from a finite admissible domain (Appendix~\ref{app:isomorphism_dictionary}, Definitions~\ref{def:bc_closure_audit}--\ref{def:rigidity_certificate_audit}).
All reported effects should be stated relative to a baseline that preserves the relevant sampling and noise structure of the measurement protocol.

\subsection{P1: rhythm ratio bias and the anti-locking index \texorpdfstring{$Q_\delta$}{Qdelta}}
\label{sec:test_P1}

\paragraph{Data.}
Cross-species or cross-condition measurements of coupled biological rhythms (e.g.\ heart--respiration, gait frequencies, circadian harmonics, neural band couplings), together with estimates of noise levels and coupling strengths.

\paragraph{Pipeline.}
\begin{enumerate}
  \item Infer an effective frequency ratio $\alpha$ using standard phase-synchrony or coupled-oscillator identification methods \cite{Strogatz1994,PikovskyRosenblumKurths2001}.
  \item Estimate an operational locking tolerance $\delta$ from noise/coupling (or from empirical locking width) using the detuning-to-ratio conversion \eqref{eq:delta_from_detuning}.
  \item Compute $Q_\delta(\alpha)$ as in Eq.~\eqref{eq:Qdelta_def} (Appendix~\ref{app:repro_code} provides a reference implementation).
  \item Compare the distribution of $Q_\delta$ (or normalized $Q_\delta$ relative to the Diophantine bound) against appropriate baselines that preserve sampling and noise structure.
\end{enumerate}
\paragraph{Decision criterion.}
If there is no anti-locking selection pressure, $Q_\delta$ should not systematically exceed baselines.
If the control-law hypothesis holds (Remark~\ref{rem:control_law}), $Q_\delta$ should be significantly shifted upward and should exhibit an enhanced upper tail near the golden-branch extremum.
\paragraph{Audit form (baseline and rejection region).}
Fix an admissible protocol domain (windowing, inference method, and surrogate family) and compute the scale-free score $Z_\delta=\sqrt{\delta}\,Q_\delta$ (Eq.~\eqref{eq:Zdelta_def}) for each inferred coupling ratio.
Let $T$ be a pre-specified summary statistic of the empirical $Z_\delta$ sample (e.g.\ median or an upper-tail quantile).
Generate a baseline distribution of $T$ using protocol-matched surrogates (e.g.\ time-shuffled or phase-randomized controls that preserve power spectra and noise levels) and reject the no-bias hypothesis at level $\alpha$ if $T$ exceeds the $(1-\alpha)$ baseline quantile.

\subsection{P2: phase-friction certificate \texorpdfstring{$E_N$}{EN} vs maintenance dissipation}
\label{sec:test_P2}

\paragraph{Claim.}
Stronger AEC should manifest as lower effective mismatch density $\sigma$ (or lower discrepancy-based proxies such as $D_N^\ast$ at matched scales) under comparable external perturbations, at the cost of higher maintenance dissipation $\dot W_{\mathrm{diss}}$.

\paragraph{Operational proxies.}
Fix a readout variable $y(t)$ (e.g.\ inter-event intervals, phase-difference increments, or any scalar readout chosen by the protocol) and map it to a phase coordinate $u(t)\in[0,1)$ by uniformization under a reference distribution:
\begin{equation}
  u(t) := \widehat{F}\bigl(y(t)\bigr),
  \label{eq:u_uniformization}
\end{equation}
where $\widehat{F}$ is an empirically estimated CDF under matched baseline conditions.
For a window of $N$ samples ending at time $t$, compute the star discrepancy $D_N^\ast(t)$ of $\{u_{t,1},\dots,u_{t,N}\}$ and define
\begin{equation}
  E_N(t):=N D_N^\ast(t),\qquad S_{\mathrm{pf}}^{(N)}(t):=\kB E_N(t).
\end{equation}
Define the external-time mismatch rate by a finite difference
\begin{equation}
  \Sigma_N(t) := \frac{E_N(t+\Delta t)-E_N(t)}{\Delta t},
  \label{eq:SigmaN_def}
\end{equation}
which estimates $\Sigma(x,t)$ after smoothing (Section~\ref{sec:mismatch_density}).

For the dissipation side, use a protocol-matched estimate of maintenance power $P_{\mathrm{maint}}(t)$ (e.g.\ metabolic power minus external mechanical work, or ATP expenditure rates). In the minimal closure where maintenance is dominated by exporting phase-friction entropy at temperature $T_c$, the bound \eqref{eq:Wdiss_lower} gives
\begin{equation}
  P_{\mathrm{maint}}(t)\ \ge\ \kB T_c\,\Sigma(x,t),
  \label{eq:Pmaint_lower}
\end{equation}
and with geometric impedance one expects an affine lower envelope
\begin{equation}
  P_{\mathrm{maint}}(t)\ \gtrsim\ \kB T_c\,\Sigma(x,t)\ +\ P_{\mathrm{geom}}(t),
  \label{eq:Pmaint_affine}
\end{equation}
where $P_{\mathrm{geom}}$ aggregates architecture-dependent overhead terms (Section~\ref{sec:geometric_landauer}).

\paragraph{Decision criterion.}
Across matched conditions, systems with enhanced repair/control should exhibit reduced $\Sigma_N$ at the relevant scale.
For quantitative fitting, one can regress $P_{\mathrm{maint}}$ against $\Sigma_N$ to estimate $\kB T_c$ as the minimal slope and diagnose additional intercept/overhead terms associated with $Z_{\mathrm{geom}}$.
To make this auditable, the window length $N$, step size $\Delta t$, and smoothing operator used to compute $\Sigma_N$ should be fixed (or selected by a bounded-complexity rule) and reported alongside the fitted slope/intercept and their uncertainty under resampling.

\subsection{P3: hierarchical relaxation and \texorpdfstring{$1/f$}{1/f} structure from Fibonacci/Zeckendorf layering}
\label{sec:test_P3}

HPA--$\Omega$ predicts that canonical multi-scale coding (Ostrowski/Zeckendorf hierarchies) can imprint hierarchical time scales on readout-induced residuals \cite{Ma2025HPT}.
If biological repair/relaxation processes aggregate across a roughly logarithmic ladder of time scales, a mid-band $1/f$ spectrum can appear as a robust template (with task- and protocol-dependent prefactors).
This yields a falsifiable interface: identify the relevant band, estimate the slope and prefactor, and test consistency with hierarchical-ladder aggregation rather than with purely white-noise assumptions.
For background on $1/f$ phenomenology and spectral estimation practices, see e.g.\ \cite{Milotti2002}.

\subsection{Reference implementations}
\label{sec:reference_impl}

Appendix~\ref{app:repro_code} includes minimal Python reference implementations for:
\begin{itemize}
  \item logarithmic mismatch-growth compatibility for the golden branch vs linear growth for rationals (Experiment~A),
  \item a numerical illustration of the predictive-information-rate threshold \eqref{eq:Idot_threshold} (Experiment~B),
  \item computation of $Q_\delta$ (Experiment~C).
\end{itemize}
The scripts are provided for transparent reproducibility of the certificate computations and threshold numerics used in the text.


